\subsection{Ficha 1: Identificación de microservicios al descomponer un monolito}
\begin{description}
  \item[Referencia (IEEE):] \fullcite{Oumoussa_2024}
  \item[DOI o ISBN verificado:] \url{https://doi.org/10.1109/ACCESS.2024.3365079}
  \item[Resumen:]El articulo aborda la transformación de sistemas monolíticos hacia arquitecturas de microservicios (MSA) ofrece ventajas como escalabilidad, agilidad y adopción en la nube. Aun así, el éxito de esta depende de una correcta y eficiente identificación de los microservicios, área que permanece aún en desarrollo. Separar un sistema grande en microservicios es difícil porque presenta varias dificultades, ya que dividir un sistema grande en servicios independientes requiere identificar correctamente sus componentes. Para tratar esta problemática, esta revisión sistemática de la literatura exploró tres preguntas clave. En total los autores seleccionaron 168 artículos para realizar este análisis. Estos hallazgos ponen en manifiesto áreas claves que requieren atención, con la necesidad de herramientas más automatizadas y precisas, el establecimiento de criterios estandarizados para su evaluación y mayor entendimiento de factores humanos que influyen en transiciones exitosas.
Las métricas de evaluación no se aplican de manera consistente en todos los estudios, por lo tanto, es difícil comparar la efectividad de los métodos entre estudios.
  \item[Dato concreto:] Los autores identificaron 5 objetivos de investigación en el campo de identificación de microservicios, y una variedad de métricas de evaluación, siendo las métricas de cohesión, IFN y LOC las más comunes. Ubicación: sección IV-B (Identificación, investigación y evaluación de microservicios), página 23393. 
  \item[Nota personal:] Me ayuda a entender que la transformación de sistemas monolíticos a arquitecturas de microservicios es un proceso que presenta varias dificultades, y que depende críticamente de aplicar métricas de evaluación precisas (como cohesión y LOC) para separar correctamente los componentes, y lo usaría para para explicar cómo está dividido el sistema que elijas en componentes o servicios, y con qué criterio se separó cada uno.
\end{description}

\subsection{Ficha 2: Descomposición de monolitos en microservicios}
\begin{description}
  \item[Referencia (IEEE):] \fullcite{Abgaz_2023}
  \item[DOI o ISBN verificado:] \url{https://doi.org/10.1109/TSE.2023.3287297}
  \item[Resumen:] Muchas empresas que utilizan sistemas complejos buscan exhaustivamente soluciones para descomponer sus aplicaciones monolíticas en lo que son microservicios. Se examinaron 35 artículos de investigación mediante una revisión sistemática de Literatura. Los autores clasifican los enfoques de descomposición según el tipo de dato en el que se basan: basado en modelos, basado en código basado en registros y basado en versiones 
Se observo una ausencia de una descripción del proceso de descomposición de extremo a extremo. Como hallazgo se observó que la descomposición de monolitos aún se encuentra en una etapa temprana, y la existencia de falta de métodos y de herramientas de apoyo.
  \item[Dato concreto:] \item[Dato concreto:] Los autores concluyen que las fases principales del M2MDF son: recopilación de datos de entrada, análisis de monolitos, identificación de microservicios, optimización, evaluación e implementación de microservicios. Ubicación: sección VIII (Conclusión), página 4236.
  \item[Nota personal:] Hay varios enfoques para descomponer, no uno solo, y que aun se encuentra con la existencia de falta de métodos y herramientas de apoyo. 
Lo usuario para poder expresar cómo un sistema real evolucionó de un monolito a servicios.
\end{description}

\subsection{Ficha 3: Prácticas de seguridad en sistemas de microservicios}
\begin{description}
  \item[Referencia (IEEE):] \fullcite{Rezaei_Nasab_2023}
  \item[DOI o ISBN verificado:] \url{https://doi.org/10.1016/j.jss.2022.111563}
  \item[Resumen:] La seguridad en los sistemas de microservicios es un gran reto, principalmente por la brecha de conocimiento que existe entre los profesionales sobre cómo asegurarlos correctamente. Para ayudar con esto, el estudio extrajo información de repositorios y foros para proponer prácticas de seguridad en áreas clave como autorización, manejo de tokens y bases de datos. Como método de validación, en lugar de implementarlas en sistemas reales, realizaron una encuesta a 74 profesionales de la industria. Como resultado, los encuestados consideraron que estas prácticas eran útiles para su trabajo
  \item[Dato concreto:] El API Gateway para sistemas grandes (PA4) fue la más aceptada: 70,27 % la consideró “absolutamente útil” y 25,68 % “útil”, o sea 95,95 % en total. Está en la Tabla 3, sección 4.2.1.
  \item[Nota personal:] Me ayuda a entender que las prácticas de seguridad se organizan en varias categorías como autorización y tokens. Lo usuario para argumentar que un sistema distribuido tiene más puntos que proteger y eso es una limitación.
\end{description}

\subsection{Ficha 4: Herramientas de monitoreo para DevOps y microservicios}
\begin{description}
  \item[Referencia (IEEE):] \fullcite{Giamattei_2024}
  \item[DOI o ISBN verificado:] \url{https://doi.org/10.1016/j.jss.2023.111906}
  \item[Resumen:]El desarrollo y operación de los sistemas basados en microservicios es una situación compleja, el principal problema que aborda es la dificultad de seleccionar una herramienta de monitorización adecuada para sistemas de microservicios desarrollados bajo practicas DevOps.
Se emplea una revisión sistemática de literatura gris, ejecutando un proceso riguroso de búsqueda, selección y extracción de datos en repositorios como GitHub y motores de búsqueda, lo que llevo a 71 herramientas para analizar.
Como resultado de analizar esas 71, contrastaron los resultados y sintetizaron el mapa de las características emergentes de las herramientas analizadas, sus ventajas y desventajas y otras características más.
  \item[Dato concreto:] La mayoría de las herramientas (58%) están dirigidas a sistemas distribuidos de alto nivel, mientras que el resto, especialmente las más nuevas, se centran específicamente en tecnologías de servicios web (21%) y en sistemas basados ​​en microservicios (21%). Ubicación: sección 8.1(Principales hallazgos y orientaciones para ingenieros DevOps), pagina 13
  \item[Nota personal:] Me ayuda a entender que monitorear es la base para detectar fallos, y que el principal problema es la dificultad de seleccionar una herramienta adecuada.
Lo usuaria para explicar como detectar que un nodo o servicio dejo de responder.
\end{description}

\subsection{Ficha 5: Visualización dinámica de arquitecturas de microservicios}
\begin{description}
  \item[Referencia (IEEE):] \fullcite{Gortney_2022}
  \item[DOI o ISBN verificado:] \url{https://doi.org/10.1109/ACCESS.2022.3221130}
  \item[Resumen:] A medida que los microservicios van evolucionando y popularizándose, se presentan diversos inconvenientes. Uno de estos es la falta de una visión completa del sistema, no permitiendo visualizar la arquitectura de una forma integral. En este articulo se hizo un estudio de mapeo investigando el análisis dinámico para visualizar la arquitectura del sistema. Durante el proceso de revisión literaria probaron con 37 consultas de búsquedas diferentes, y la consulta final fue lo suficientemente amplia como para abarcar múltiples métodos de análisis dinámico y demostró arrojar los resultados más relevantes
Los autores identificaron tres técnicas principales de análisis dinámico: minería de procesos, rastreo distribuido y monitorización, además encontraron herramientas para recopilar datos de tiempo de ejecución y crear visualización a partir de su análisis.
  \item[Dato concreto:] Existen tres técnicas principales de análisis dinámico: minería de procesos, rastreo distribuido y monitorización. Ubicación: sección 9(Conclusión), pagina 120011
  \item[Nota personal:] Me ayuda a entender que ver la arquitectura de un sistema mientras se ejecuta, con tres técnicas: minería de procesos, rastreo distribuido y monitorización.
\end{description}
